\chapter{Introduction}
\label{ch:introduction}

\section{Motivation}

Adaptive control addresses the fundamental challenge of controlling systems with uncertain or time-varying parameters. Unlike robust control, which attempts to handle uncertainty through conservative design margins, adaptive control \textit{learns} the unknown parameters online, potentially achieving superior performance while maintaining stability guarantees.

The need for adaptive control arises in numerous applications:
\begin{itemize}
    \item \textbf{Aerospace:} Aircraft dynamics change with fuel consumption, payload, and atmospheric conditions.
    \item \textbf{Robotics:} Manipulators handle objects with unknown mass and inertia.
    \item \textbf{Biomedical Systems:} Patient-specific physiological parameters vary widely and evolve over time.
    \item \textbf{Automotive:} Vehicle dynamics depend on load, tire wear, and road conditions.
    \item \textbf{Process Control:} Chemical plants exhibit parameter drift due to catalyst aging and feedstock variations.
\end{itemize}

\subsection{The Role of Memory in Adaptive Systems}

Classical adaptive control theory, developed primarily in the 1970s--1990s by pioneers such as Narendra, Åström, and Slotine, focused on memoryless (Markovian) systems where the current state fully determines future evolution. However, many physical systems exhibit \textbf{memory effects}:
\begin{itemize}
    \item \textbf{Biological Systems:} Neural plasticity, muscle fatigue, and cardiovascular dynamics involve distributed time delays and history-dependent responses.
    \item \textbf{Viscoelastic Materials:} Stress-strain relationships depend on deformation history through memory kernels.
    \item \textbf{Financial Systems:} Market dynamics incorporate memory of past prices and trading volumes.
    \item \textbf{Climate Models:} Atmospheric and oceanic processes exhibit long-range temporal correlations.
\end{itemize}

Memory effects are mathematically captured by integro-differential equations of the form:
\begin{equation}
    \dot{x}(t) = f(x(t), u(t)) + \int_{-\tau}^0 K(\xi) g(x(t+\xi)) \, d\xi,
\end{equation}
where the kernel $K(\cdot)$ weights past states $x(t+\xi)$ for $\xi \in [-\tau, 0]$.

\textbf{Challenge:} Standard Lyapunov stability theory does not directly apply to systems with memory. We require \textit{Lyapunov-Krasovskii} functionals that account for the entire state history over $[t-\tau, t]$.

\subsection{Contributions of This Thesis}

This thesis makes the following contributions to adaptive control theory:

\begin{enumerate}
    \item \textbf{Unified Mathematical Framework:} We develop a complete functional-analytic foundation incorporating:
    \begin{itemize}
        \item Sobolev spaces $W^{k,p}$ for state regularity,
        \item Barbalat's Lemma with full proof from first principles,
        \item Persistent excitation theory with spectral characterization.
    \end{itemize}

    \item \textbf{Lyapunov-Krasovskii Stability Analysis:} We construct novel Lyapunov functionals for adaptive systems with memory, providing:
    \begin{itemize}
        \item Complete algebraic derivation of time derivatives (every step shown),
        \item Proof of uniform ultimate boundedness (UUB) under bounded disturbances,
        \item Proof of global asymptotic stability (GAS) under persistent excitation.
    \end{itemize}

    \item \textbf{Formal Adaptive Update Laws:} We derive projection-based and saturation-based parameter update laws with:
    \begin{itemize}
        \item Convergence guarantees,
        \item Robustness to measurement noise,
        \item Computational efficiency.
    \end{itemize}

    \item \textbf{Robust Extensions for Practical Implementation:} We address:
    \begin{itemize}
        \item Non-PE scenarios (tracking without parameter convergence),
        \item State-dependent disturbances,
        \item Unmodeled dynamics,
        \item Adaptive safety mechanisms for explainable AI (xAI) applications.
    \end{itemize}

    \item \textbf{Physiological Application:} We apply the framework to \textit{heart-brain coupling models}, demonstrating:
    \begin{itemize}
        \item Adaptive control of cardiac rhythm via neural stimulation,
        \item Brain-computer interface (BCI) integration with OpenBCI and Neuralink-style hardware,
        \item Parameter identification in physiological oscillators (FitzHugh-Nagumo, Van der Pol).
    \end{itemize}

    \item \textbf{Complete Computational Implementation:} We provide:
    \begin{itemize}
        \item Python reference implementation with full documentation,
        \item Simulation results validating theoretical predictions,
        \item Comparison of PE versus non-PE scenarios,
        \item Practical implementation checklists.
    \end{itemize}
\end{enumerate}

\section{Historical Context}

\subsection{Classical Adaptive Control (1960s--1980s)}

The foundations of adaptive control were laid in the 1960s--1970s:
\begin{itemize}
    \item \textbf{Model Reference Adaptive Control (MRAC):} Whitaker (1958), Parks (1966), and Monopoli (1974) developed MRAC for systems tracking a reference model.
    \item \textbf{Self-Tuning Regulators:} Åström and Wittenmark (1973) introduced self-tuning regulators based on recursive parameter estimation.
    \item \textbf{Stability Theory:} Narendra and colleagues established Lyapunov-based stability proofs in the 1970s--1980s \cite{Narendra1989}.
\end{itemize}

\textbf{Key Insight:} The certainty equivalence principle—use parameter estimates $\thetahat(t)$ as if they were true parameters—can be stabilized through appropriately designed adaptive laws.

\subsection{Persistent Excitation (1970s--1980s)}

A crucial discovery was that tracking error convergence ($e \to 0$) does not guarantee parameter convergence ($\thetahat \to \thetastar$). Parameter convergence requires \textbf{persistent excitation}:
\begin{itemize}
    \item Anderson (1977) formalized PE conditions.
    \item Narendra and Annaswamy (1987) provided detailed PE analysis for MRAC.
    \item Boyd and Sastry (1986) gave spectral characterization of PE signals.
\end{itemize}

\subsection{Robust Adaptive Control (1980s--1990s)}

Practical implementation revealed challenges:
\begin{itemize}
    \item \textbf{Unmodeled Dynamics:} Real systems deviate from nominal models.
    \item \textbf{Disturbances:} External perturbations and measurement noise.
    \item \textbf{Computational Constraints:} Real-time parameter adaptation.
\end{itemize}

Robust adaptive control, developed by Ioannou and Sun (1996), Slotine and Li (1991), and Krstić, Kanellakopoulos, and Kokotović (1995), introduced:
\begin{itemize}
    \item $\sigma$-modification and $e$-modification to prevent parameter drift,
    \item Projection algorithms to enforce parameter bounds,
    \item Deadzone modifications for robustness to disturbances.
\end{itemize}

\subsection{Adaptive Control with Memory (1990s--Present)}

Interest in non-Markovian systems led to:
\begin{itemize}
    \item \textbf{Fractional-Order Systems:} Fractional calculus captures memory through non-integer derivatives.
    \item \textbf{Time-Delay Systems:} Explicit time delays in feedback (Niculescu 2001, Gu et al. 2003).
    \item \textbf{Integro-Differential Systems:} Memory kernels in viscoelasticity, population dynamics, and epidemiology.
\end{itemize}

\textbf{Gap in Literature:} Most work treats memory as a perturbation or focuses on specific kernel forms (exponential, polynomial). A \textit{unified framework} for general memory kernels with adaptive parameter uncertainty is lacking.

\textbf{This thesis fills this gap.}

\section{Physiological Application: Heart-Brain Coupling}

A motivating application is the \textbf{heart-brain coupling model (HBCM)}, which describes bidirectional interactions between cardiac and neural oscillators:
\begin{align}
    \text{Neural (FitzHugh-Nagumo):} \quad &
    \begin{cases}
        \dot{v} = v - v^3/3 - w + I_{\text{cardiac} \to \text{neural}} \\
        \dot{w} = (v + a - bw) / c
    \end{cases} \\[0.5em]
    \text{Cardiac (Van der Pol):} \quad &
    \begin{cases}
        \dot{x} = y \\
        \dot{y} = -\omega^2 x + \mu (1 - x^2) y + I_{\text{neural} \to \text{cardiac}}
    \end{cases}
\end{align}

The coupling terms involve delayed feedback:
\begin{align}
    I_{\text{cardiac} \to \text{neural}}(t) &= g_1 x(t - \tau_1), \\
    I_{\text{neural} \to \text{cardiac}}(t) &= g_2 v(t - \tau_2),
\end{align}
where $g_1, g_2$ are coupling gains and $\tau_1, \tau_2$ are propagation delays (typically 50--200 ms).

\subsection{Control Objective}

\textbf{Goal:} Regulate cardiac rhythm (heart rate, heart rate variability) via neural stimulation, adapting to unknown patient-specific parameters $(a, b, c, \mu, \omega, g_1, g_2, \tau_1, \tau_2)$.

\textbf{Challenges:}
\begin{itemize}
    \item \textbf{Parameter Uncertainty:} Physiological parameters vary across individuals and over time.
    \item \textbf{Delays:} Axonal conduction and synaptic delays introduce memory effects.
    \item \textbf{Safety:} Neural stimulation must remain within safe bounds.
    \item \textbf{Real-Time Adaptation:} Parameters must be estimated from online measurements (EEG, ECG).
\end{itemize}

This thesis develops adaptive controllers for HBCM, integrating:
\begin{itemize}
    \item Brain-computer interfaces (OpenBCI, Neuralink) for neural state measurement,
    \item Adaptive Lyapunov-based control for cardiac regulation,
    \item Safety constraints via projection and saturation.
\end{itemize}

\section{Thesis Organization}

The remainder of this thesis is organized as follows:

\begin{description}
    \item[Chapter 2: Literature Review] Comprehensive survey of adaptive control, Lyapunov stability theory, memory kernel modeling, and physiological applications.

    \item[Chapter 3: Mathematical Preliminaries] Function spaces ($\Lp{p}$, $W^{k,p}$), Barbalat's Lemma with complete proof, persistent excitation theory.

    \item[Chapter 4: System Definition and Augmented States] Formal problem statement, uncertainty modeling, augmented memory state construction.

    \item[Chapter 5: Lyapunov-Krasovskii Stability Analysis] Construction of Lyapunov functional, complete derivative algebra, stability conditions.

    \item[Chapter 6: Adaptive Update Laws] Derivation of gradient-based, projection-based, and saturation-based adaptive laws with convergence analysis.

    \item[Chapter 7: Main Stability Theorems] Formal statements and proofs of GAS under PE, UUB with disturbances, convergence rates.

    \item[Chapter 8: Robust Extensions] Non-PE scenarios, unmodeled dynamics, disturbance rejection, adaptive safety for xAI.

    \item[Chapter 9: Simulation Results] Python/MATLAB implementations, numerical experiments, PE vs. non-PE comparisons, parameter convergence plots.

    \item[Chapter 10: Practical Guidelines] Implementation checklists, tuning heuristics, convergence diagnostics, safety verification.

    \item[Chapter 11: Conclusion] Summary of contributions, future research directions.

    \item[Appendices] Detailed derivations, additional lemmas, complete code listings, extended results.
\end{description}

\section{Notation and Conventions}

Throughout this thesis, we use the following notation:
\begin{itemize}
    \item $\R, \C, \N$: Real numbers, complex numbers, natural numbers
    \item $\R^n$: $n$-dimensional Euclidean space
    \item $\norm{\cdot}$: Euclidean norm for vectors, induced 2-norm for matrices
    \item $\Lp{p}[0,\infty)$: Lebesgue space with norm $\norm{f}_{\Lp{p}} = (\int_0^\infty |f(t)|^p \, dt)^{1/p}$
    \item $W^{k,p}$: Sobolev space of functions with $k$ weak derivatives in $\Lp{p}$
    \item $A \succ 0$: Matrix $A$ is positive definite
    \item $\lambda_{\min}(A), \lambda_{\max}(A)$: Minimum and maximum eigenvalues of $A$
    \item $\thetastar$: True parameter vector
    \item $\thetahat$: Parameter estimate
    \item $\thetatilde = \thetahat - \thetastar$: Parameter estimation error
    \item $e(t) = x(t) - x_r(t)$: Tracking error
    \item PE: Persistently exciting
    \item GAS: Globally asymptotically stable
    \item UUB: Uniformly ultimately bounded
\end{itemize}

Matrix inequalities are in the positive (semi)definite sense: $A \succeq B$ means $A - B$ is positive semidefinite.

\section{Contributions to Knowledge}

This thesis advances the state of knowledge in adaptive control through:
\begin{enumerate}
    \item \textbf{Theoretical Rigor:} Complete, self-contained proofs of all results with every algebraic step shown.
    \item \textbf{Generality:} Unified treatment of memory kernels, not restricted to specific forms.
    \item \textbf{Practical Applicability:} Implementation guidelines, safety mechanisms, and real-world application (HBCM).
    \item \textbf{Reproducibility:} Full open-source code with documentation, enabling verification and extension.
\end{enumerate}

The work is positioned at the intersection of control theory, functional analysis, and biomedical engineering, contributing to all three fields.
